% Original diagrams; all four panels show the unconditioned graph.
\begin{tikzpicture}[x=1mm,y=1mm,font=\small,
  pgm variable/.style={circle,draw=BookInk,fill=BookPaper,
    minimum size=7mm,inner sep=0pt},
  middle/.style={pgm variable,draw=BookTeal,fill=BookTeal!9!white},
  edge/.style={-{Stealth[length=2mm]},draw=BookInk,thick},
  heading/.style={font=\heiti\footnotesize,text=BookInk}]
  \node[heading] at (26,60) {直接连接};
  \node[pgm variable] (dx) at (10,47) {$X$};
  \node[pgm variable] (dy) at (42,47) {$Y$};
  \draw[edge] (dx)--(dy);
  \node[font=\footnotesize,text=BookMuted] at (26,37)
    {无内部阻断点};

  \node[heading] at (82,60) {链式关系};
  \node[pgm variable] (cx) at (64,47) {$X$};
  \node[middle] (cz) at (82,47) {$Z$};
  \node[pgm variable] (cy) at (100,47) {$Y$};
  \draw[edge] (cx)--(cz);
  \draw[edge] (cz)--(cy);
  \node[font=\footnotesize,text=BookMuted] at (82,37)
    {给定$Z$后阻断};

  \node[heading] at (26,25) {共同原因：分叉};
  \node[pgm variable] (fx) at (8,3) {$X$};
  \node[middle] (fz) at (26,14) {$Z$};
  \node[pgm variable] (fy) at (44,3) {$Y$};
  \draw[edge] (fz)--(fx);
  \draw[edge] (fz)--(fy);
  \node[font=\footnotesize,text=BookMuted] at (26,-8)
    {给定$Z$后阻断};

  \node[heading] at (82,25) {共同结果：碰撞};
  \node[pgm variable] (vx) at (64,14) {$X$};
  \node[middle] (vz) at (82,3) {$Z$};
  \node[pgm variable] (vy) at (100,14) {$Y$};
  \draw[edge] (vx)--(vz);
  \draw[edge] (vy)--(vz);
  \node[font=\footnotesize,text=BookMuted] at (82,-8)
    {给定$Z$后可能打开};
\end{tikzpicture}
